% Options for packages loaded elsewhere
% Options for packages loaded elsewhere
\PassOptionsToPackage{unicode}{hyperref}
\PassOptionsToPackage{hyphens}{url}
\PassOptionsToPackage{dvipsnames,svgnames,x11names}{xcolor}
%
\documentclass[
  11pt,
  letterpaper,
  DIV=11,
  numbers=noendperiod]{scrartcl}
\usepackage{xcolor}
\usepackage[left=2.5cm,right=2.5cm,top=2.5cm,bottom=2.5cm]{geometry}
\usepackage{amsmath,amssymb}
\setcounter{secnumdepth}{5}
\usepackage{iftex}
\ifPDFTeX
  \usepackage[T1]{fontenc}
  \usepackage[utf8]{inputenc}
  \usepackage{textcomp} % provide euro and other symbols
\else % if luatex or xetex
  \usepackage{unicode-math} % this also loads fontspec
  \defaultfontfeatures{Scale=MatchLowercase}
  \defaultfontfeatures[\rmfamily]{Ligatures=TeX,Scale=1}
\fi
\usepackage{lmodern}
\ifPDFTeX\else
  % xetex/luatex font selection
  \setmainfont[]{Times New Roman}
\fi
% Use upquote if available, for straight quotes in verbatim environments
\IfFileExists{upquote.sty}{\usepackage{upquote}}{}
\IfFileExists{microtype.sty}{% use microtype if available
  \usepackage[]{microtype}
  \UseMicrotypeSet[protrusion]{basicmath} % disable protrusion for tt fonts
}{}
\makeatletter
\@ifundefined{KOMAClassName}{% if non-KOMA class
  \IfFileExists{parskip.sty}{%
    \usepackage{parskip}
  }{% else
    \setlength{\parindent}{0pt}
    \setlength{\parskip}{6pt plus 2pt minus 1pt}}
}{% if KOMA class
  \KOMAoptions{parskip=half}}
\makeatother
% Make \paragraph and \subparagraph free-standing
\makeatletter
\ifx\paragraph\undefined\else
  \let\oldparagraph\paragraph
  \renewcommand{\paragraph}{
    \@ifstar
      \xxxParagraphStar
      \xxxParagraphNoStar
  }
  \newcommand{\xxxParagraphStar}[1]{\oldparagraph*{#1}\mbox{}}
  \newcommand{\xxxParagraphNoStar}[1]{\oldparagraph{#1}\mbox{}}
\fi
\ifx\subparagraph\undefined\else
  \let\oldsubparagraph\subparagraph
  \renewcommand{\subparagraph}{
    \@ifstar
      \xxxSubParagraphStar
      \xxxSubParagraphNoStar
  }
  \newcommand{\xxxSubParagraphStar}[1]{\oldsubparagraph*{#1}\mbox{}}
  \newcommand{\xxxSubParagraphNoStar}[1]{\oldsubparagraph{#1}\mbox{}}
\fi
\makeatother


\usepackage{longtable,booktabs,array}
\newcounter{none} % for unnumbered tables
\usepackage{calc} % for calculating minipage widths
% Correct order of tables after \paragraph or \subparagraph
\usepackage{etoolbox}
\makeatletter
\patchcmd\longtable{\par}{\if@noskipsec\mbox{}\fi\par}{}{}
\makeatother
% Allow footnotes in longtable head/foot
\IfFileExists{footnotehyper.sty}{\usepackage{footnotehyper}}{\usepackage{footnote}}
\makesavenoteenv{longtable}
\usepackage{graphicx}
\makeatletter
\newsavebox\pandoc@box
\newcommand*\pandocbounded[1]{% scales image to fit in text height/width
  \sbox\pandoc@box{#1}%
  \Gscale@div\@tempa{\textheight}{\dimexpr\ht\pandoc@box+\dp\pandoc@box\relax}%
  \Gscale@div\@tempb{\linewidth}{\wd\pandoc@box}%
  \ifdim\@tempb\p@<\@tempa\p@\let\@tempa\@tempb\fi% select the smaller of both
  \ifdim\@tempa\p@<\p@\scalebox{\@tempa}{\usebox\pandoc@box}%
  \else\usebox{\pandoc@box}%
  \fi%
}
% Set default figure placement to htbp
\def\fps@figure{htbp}
\makeatother





\setlength{\emergencystretch}{3em} % prevent overfull lines

\providecommand{\tightlist}{%
  \setlength{\itemsep}{0pt}\setlength{\parskip}{0pt}}



 


\usepackage{booktabs}
\usepackage{longtable}
\usepackage{array}
\usepackage{multirow}
\usepackage{wrapfig}
\usepackage{float}
\usepackage{colortbl}
\usepackage{pdflscape}
\usepackage{tabu}
\usepackage{threeparttable}
\usepackage{threeparttablex}
\usepackage[normalem]{ulem}
\usepackage{makecell}
\usepackage{xcolor}
\usepackage{float}
\usepackage{placeins}
\usepackage{xcolor}
\usepackage[left]{lineno}
\modulolinenumbers[1]
\KOMAoption{captions}{tableheading}
\makeatletter
\@ifpackageloaded{caption}{}{\usepackage{caption}}
\AtBeginDocument{%
\ifdefined\contentsname
  \renewcommand*\contentsname{Table of contents}
\else
  \newcommand\contentsname{Table of contents}
\fi
\ifdefined\listfigurename
  \renewcommand*\listfigurename{List of Figures}
\else
  \newcommand\listfigurename{List of Figures}
\fi
\ifdefined\listtablename
  \renewcommand*\listtablename{List of Tables}
\else
  \newcommand\listtablename{List of Tables}
\fi
\ifdefined\figurename
  \renewcommand*\figurename{Figure}
\else
  \newcommand\figurename{Figure}
\fi
\ifdefined\tablename
  \renewcommand*\tablename{Table}
\else
  \newcommand\tablename{Table}
\fi
}
\@ifpackageloaded{float}{}{\usepackage{float}}
\floatstyle{ruled}
\@ifundefined{c@chapter}{\newfloat{codelisting}{h}{lop}}{\newfloat{codelisting}{h}{lop}[chapter]}
\floatname{codelisting}{Listing}
\newcommand*\listoflistings{\listof{codelisting}{List of Listings}}
\makeatother
\makeatletter
\makeatother
\makeatletter
\@ifpackageloaded{caption}{}{\usepackage{caption}}
\@ifpackageloaded{subcaption}{}{\usepackage{subcaption}}
\makeatother
\usepackage{bookmark}
\IfFileExists{xurl.sty}{\usepackage{xurl}}{} % add URL line breaks if available
\urlstyle{same}
\hypersetup{
  pdftitle={From One-Cluster SPC to Multicluster SPC: Distribution-Free Control Charts for Cleaning Validation Stability Analysis},
  colorlinks=true,
  linkcolor={blue},
  filecolor={Maroon},
  citecolor={Blue},
  urlcolor={Blue},
  pdfcreator={LaTeX via pandoc}}


\title{From One-Cluster SPC to Multicluster SPC: Distribution-Free
Control Charts for Cleaning Validation Stability Analysis}
\author{}
\date{}
\begin{document}
\maketitle


\floatplacement{table}{H}
\floatplacement{figure}{H}

\section{\texorpdfstring{\texttt{d\_ccssbu\_l} Function and
\texttt{d\_ccssbu\_l\_d\_selection\_arl1}
Function}{d\_ccssbu\_l Function and d\_ccssbu\_l\_d\_selection\_arl1 Function}}\label{d_ccssbu_l-function-and-d_ccssbu_l_d_selection_arl1-function}

\begin{itemize}
\tightlist
\item
  USL is used to form natural clusters in code but we use API to form
  natural clusters in the paper theoretical part.
\item
  Functions\texttt{d\_ccssbu\_l()} can automatically detect 3A or 3B
  stages to generate control charts, algorithm's out in table.
\item
  If we set \texttt{n\_stage3A\_events} to be a very large number, then
  the analysis will essentially be retrospective analysis.
\item
  \texttt{d\_ccssbu\_l\_d\_selection\_arl1()} used
  \texttt{d\_ccssbu\_l()} at \texttt{fit\_q} and \texttt{fit\_i} part.
\end{itemize}

\subsection{\texorpdfstring{d-CCSSBU-L Selection of \(d\) and ARL1
Results for Equipment
A}{d-CCSSBU-L Selection of d and ARL1 Results for Equipment A}}\label{d-ccssbu-l-selection-of-d-and-arl1-results-for-equipment-a}

\begin{itemize}
\tightlist
\item
  run the function \texttt{d\_ccssbu\_l\_d\_selection\_arl1()}
\item
  generate the csv file
  \texttt{d\_CCSSBU\_L\_ARL1\_Equipment\_A\_tbl\_wide.csv} use
  \texttt{readr::write\_csv()}
\item
  generate Table 1 in the paper by the function
  \texttt{render\_d\_ccssbu\_l\_recommendation\_tbl()} using argument
  \texttt{tbl\ =\ \ d\_CCSSBU\_L\_ARL1\_Equipment\_A\_recommendation}
\item
  generate ARL\(_1\) curve by
  \texttt{plot\_d\_ccssbu\_l\_arl1\_curves()} and the curve picture name
  is \texttt{d-ccssbu-l-arl1-curves-equipment-a.jpeg}
\end{itemize}

\begin{table}[!h]
\centering
\caption{\label{tab:d-ccssbu-l-recommendation-equipment-a}Recommended d values for the d-CCSSBU-L chart for Equipment A.}
\centering
\resizebox{\ifdim\width>\linewidth\linewidth\else\width\fi}{!}{
\fontsize{8}{10}\selectfont
\begin{tabular}[t]{lccclcllccl}
\toprule
Residue & Recommended d & q & Decision case & Observed-behavior case & ARL1 case & A & U & W(d) & G(d) & Selection rule\\
\midrule
DAR & 3 & 3 & Case 1 & D1 intersection D2 & U & \{1, 2, 3\} & \{1, 2, 3\} & 1.0000 & 0.5659 & U case: maximum d in U according to Rule 3\\
CAR & 1 & 1 & Case 2 & D1 intersection D2 & XO & \{1\} & Empty & 36.0501 & 9.7164 & XO case: minimum W over A; minimum G among candidates tied on W; maximum d if a tie remains according to Rule 3\\
Mic & 1 & 1 & Case 2 & D1 intersection D2 & XO & \{1\} & Empty & 4.3094 & 0.8249 & XO case: minimum W over A; minimum G among candidates tied on W; maximum d if a tie remains according to Rule 3\\
\bottomrule
\end{tabular}}
\end{table}

\begin{figure}

\centering{

\includegraphics[width=0.7\linewidth,height=0.5\textheight]{d-ccssbu-l-arl1-curves-equipment-a.jpeg}

}

\caption{\label{fig-d-selection-curve-a-l}Candidate-specific d-CCSSBU-L
ARL1 curves for Equipment A.}

\end{figure}%

\subsection{\texorpdfstring{d-CCSSBU-L Selection of \(d\) and ARL1
Results for Equipment
B}{d-CCSSBU-L Selection of d and ARL1 Results for Equipment B}}\label{d-ccssbu-l-selection-of-d-and-arl1-results-for-equipment-b}

Repeat the procedure of above Equipment A

\begin{table}[!h]
\centering
\caption{\label{tab:d-ccssbu-l-recommendation-equipment-b}Recommended d values for the d-CCSSBU-L chart for Equipment B.}
\centering
\resizebox{\ifdim\width>\linewidth\linewidth\else\width\fi}{!}{
\fontsize{8}{10}\selectfont
\begin{tabular}[t]{lccclcllccl}
\toprule
Residue & Recommended d & q & Decision case & Observed-behavior case & ARL1 case & A & U & W(d) & G(d) & Selection rule\\
\midrule
DAR & 2 & 26 & Case 1 & D1 intersection D2 & U & \{1, 2, 3, 4, 5, 6, 7, 8, 9\} & \{2\} & 0.9999 & 0.4141 & U case: maximum d in U according to Rule 3\\
CAR & 2 & 2 & Case 2 & D1 intersection D2 & XO & \{2\} & Empty & 8.0175 & 2.2513 & XO case: minimum W over A; minimum G among candidates tied on W; maximum d if a tie remains according to Rule 3\\
Mic & 1 & 1 & Case 2 & D1 intersection D2 & XO & \{1\} & Empty & 26.0636 & 4.3581 & XO case: minimum W over A; minimum G among candidates tied on W; maximum d if a tie remains according to Rule 3\\
\bottomrule
\end{tabular}}
\end{table}

\begin{figure}

\centering{

\includegraphics[width=0.7\linewidth,height=0.5\textheight]{d-ccssbu-l-arl1-curves-equipment-b.jpeg}

}

\caption{\label{fig-d-selection-curve-b-l}Candidate-specific d-CCSSBU-L
ARL1 curves for Equipment B.}

\end{figure}%

Above is for location chart and its ARL\(_1\) curve and d selection. The
following is for scale chart

\section{\texorpdfstring{\texttt{d\_ccssbu\_s} Function and
\texttt{d\_ccssbu\_s\_d\_selection\_arl1}
Function}{d\_ccssbu\_s Function and d\_ccssbu\_s\_d\_selection\_arl1 Function}}\label{d_ccssbu_s-function-and-d_ccssbu_s_d_selection_arl1-function}

\begin{itemize}
\tightlist
\item
  USL is used to form natural clusters in code but we use API to form
  natural clusters in the paper theoretical part.
\item
  Functions\texttt{d\_ccssbu\_s()} can automatically detect 3A or 3B
  stages to generate control charts, algorithm's out in table.
\item
  If we set \texttt{n\_stage3A\_events} to be a very large number, then
  the analysis will essentially be retrospective analysis.
\item
  \texttt{d\_ccssbu\_s\_d\_selection\_arl1()} used
  \texttt{d\_ccssbu\_s()} at \texttt{fit\_q} and \texttt{fit\_i} part.
\end{itemize}

\subsection{\texorpdfstring{d-CCSSBU-S Selection of \(d\) and ARL1
Results for Equipment
A}{d-CCSSBU-S Selection of d and ARL1 Results for Equipment A}}\label{d-ccssbu-s-selection-of-d-and-arl1-results-for-equipment-a}

\begin{itemize}
\tightlist
\item
  run the function \texttt{d\_ccssbu\_s\_d\_selection\_arl1()}
\item
  generate the csv file
  \texttt{d\_CCSSBU\_s\_ARL1\_Equipment\_A\_tbl\_wide.csv} use
  \texttt{readr::write\_csv()}
\item
  generate Table 3 in the paper by the function
  \texttt{render\_d\_ccssbu\_s\_recommendation\_tbl()} using argument
  \texttt{tbl\ =\ \ d\_CCSSBU\_S\_ARL1\_Equipment\_A\_recommendation}
\item
  generate ARL\(_1\) curve by
  \texttt{plot\_d\_ccssbu\_s\_arl1\_curves()} and the curve picture name
  is \texttt{d-ccssbu-s-arl1-curves-equipment-a.jpeg}
\end{itemize}

\begin{table}[!h]
\centering
\caption{\label{tab:d-ccssbu-s-recommendation-equipment-a}Recommended d values for the d-CCSSBU-S chart for Equipment A.}
\centering
\resizebox{\ifdim\width>\linewidth\linewidth\else\width\fi}{!}{
\fontsize{7}{9}\selectfont
\begin{tabular}[t]{lrrlllllrrl}
\toprule
Residue & d & q & Case & Admissible Set & ARL Case & A & U & W & G & Rule\\
\midrule
DAR & 1 & 3 & Case 2 & D1 ∩ D2 & XO & 1, 2 &  & 4.4881 & 1.5006 & XO case: minimum W over A; minimum G among candidates tied on W; maximum d if a tie remains according to Rule 3\\
CAR & 1 & 1 & Case 2 & D1 ∩ D2 & XO & 1 &  & 4.5574 & 2.6012 & XO case: minimum W over A; minimum G among candidates tied on W; maximum d if a tie remains according to Rule 3\\
Mic & 1 & 1 & Case 2 & D1 ∩ D2 & XO & 1 &  & 1.5395 & 0.1828 & XO case: minimum W over A; minimum G among candidates tied on W; maximum d if a tie remains according to Rule 3\\
\bottomrule
\end{tabular}}
\end{table}

\begin{figure}

\centering{

\includegraphics[width=0.7\linewidth,height=0.5\textheight]{d-ccssbu-s-arl1-curves-equipment-a.jpeg}

}

\caption{\label{fig-d-selection-curve-a-s}Candidate-specific d-CCSSBU-S
ARL1 curves for Equipment A.}

\end{figure}%

\subsection{\texorpdfstring{d-CCSSBU-S Selection of \(d\) and ARL\(_1\)
Results for Equipment
B}{d-CCSSBU-S Selection of d and ARL\_1 Results for Equipment B}}\label{d-ccssbu-s-selection-of-d-and-arl_1-results-for-equipment-b}

Repeat as above Equipment A

\begin{table}[!h]
\centering
\caption{\label{tab:d-ccssbu-s-recommendation-equipment-b}Recommended d values for the d-CCSSBU-S chart for Equipment B.}
\centering
\resizebox{\ifdim\width>\linewidth\linewidth\else\width\fi}{!}{
\fontsize{7}{9}\selectfont
\begin{tabular}[t]{lrrlllllrrl}
\toprule
Residue & d & q & Case & Admissible Set & ARL Case & A & U & W & G & Rule\\
\midrule
DAR & 13 & 26 & Case 4 & D1 & XO & 1, 2, 3, 4, 5, 6, 7, 8, 9, 10, 11, 12, 13, 14, 15, 16, 17, 18, 19, 20, 21, 22, 23, 24, 25, 26 &  & 7.3323 & 3.0336 & XO case: minimum W over A; minimum G among candidates tied on W; maximum d if a tie remains according to Rule 3\\
CAR & 1 & 2 & Case 2 & D1 ∩ D2 & XO & 1, 2 &  & 1.5221 & 1.1434 & XO case: minimum W over A; minimum G among candidates tied on W; maximum d if a tie remains according to Rule 3\\
Mic & 1 & 1 & Case 2 & D1 ∩ D2 & XO & 1 &  & 35.9163 & 3.4234 & XO case: minimum W over A; minimum G among candidates tied on W; maximum d if a tie remains according to Rule 3\\
\bottomrule
\end{tabular}}
\end{table}

\begin{figure}

\centering{

\includegraphics[width=0.7\linewidth,height=0.5\textheight]{d-ccssbu-s-arl1-curves-equipment-b.jpeg}

}

\caption{\label{fig-d-selection-curve-b-s}Candidate-specific d-CCSSBU-S
ARL1 curves for Equipment B.}

\end{figure}%

\section{\texorpdfstring{d-CCSSBU-L Charts for Recommended \(d\) Using
Function
\texttt{d\_ccssbu\_l()}}{d-CCSSBU-L Charts for Recommended d Using Function d\_ccssbu\_l()}}\label{d-ccssbu-l-charts-for-recommended-d-using-function-d_ccssbu_l}

\pandocbounded{\includegraphics[keepaspectratio]{multicluster_SPC_paper_code_files/figure-pdf/unnamed-chunk-1-1.pdf}}

\pandocbounded{\includegraphics[keepaspectratio]{multicluster_SPC_paper_code_files/figure-pdf/unnamed-chunk-1-2.pdf}}

\pandocbounded{\includegraphics[keepaspectratio]{multicluster_SPC_paper_code_files/figure-pdf/unnamed-chunk-1-3.pdf}}

\pandocbounded{\includegraphics[keepaspectratio]{multicluster_SPC_paper_code_files/figure-pdf/unnamed-chunk-1-4.pdf}}

\pandocbounded{\includegraphics[keepaspectratio]{multicluster_SPC_paper_code_files/figure-pdf/unnamed-chunk-1-5.pdf}}

\pandocbounded{\includegraphics[keepaspectratio]{multicluster_SPC_paper_code_files/figure-pdf/unnamed-chunk-1-6.pdf}}

\section{\texorpdfstring{d-CCSSBU-S Charts for Recommended \(d\) Using
Function
\texttt{d\_ccssbu\_s()}}{d-CCSSBU-S Charts for Recommended d Using Function d\_ccssbu\_s()}}\label{d-ccssbu-s-charts-for-recommended-d-using-function-d_ccssbu_s}

\pandocbounded{\includegraphics[keepaspectratio]{multicluster_SPC_paper_code_files/figure-pdf/generate-d-ccssbu-s-charts-1.pdf}}

\pandocbounded{\includegraphics[keepaspectratio]{multicluster_SPC_paper_code_files/figure-pdf/generate-d-ccssbu-s-charts-2.pdf}}

\pandocbounded{\includegraphics[keepaspectratio]{multicluster_SPC_paper_code_files/figure-pdf/generate-d-ccssbu-s-charts-3.pdf}}

\pandocbounded{\includegraphics[keepaspectratio]{multicluster_SPC_paper_code_files/figure-pdf/generate-d-ccssbu-s-charts-4.pdf}}

\pandocbounded{\includegraphics[keepaspectratio]{multicluster_SPC_paper_code_files/figure-pdf/generate-d-ccssbu-s-charts-5.pdf}}

\pandocbounded{\includegraphics[keepaspectratio]{multicluster_SPC_paper_code_files/figure-pdf/generate-d-ccssbu-s-charts-6.pdf}}

\section{d-CCSSBU Charts for Equipment A and B Using Generated jpeg
Files}\label{d-ccssbu-charts-for-equipment-a-and-b-using-generated-jpeg-files}

\subsection{d-CCSSBU Charts for DAR of Equipment A and B Using Generated
jpeg
Files}\label{d-ccssbu-charts-for-dar-of-equipment-a-and-b-using-generated-jpeg-files}

\begin{figure}[H]
\centering
\begin{tabular}{cc}
\includegraphics[width=0.48\textwidth]{d_CCSSBU_L_DAR_Equipment_A.jpeg} &
\includegraphics[width=0.48\textwidth]{d_CCSSBU_L_DAR_Equipment_B.jpeg} \\

\includegraphics[width=0.48\textwidth]{d_CCSSBU_S_DAR_Equipment_A.jpeg} &
\includegraphics[width=0.48\textwidth]{d_CCSSBU_S_DAR_Equipment_B.jpeg} \\
\end{tabular}
\caption{d-CCSSBU-L and d-CCSSBU-S charts for DAR of Equipment A and B.}
\label{fig:d-ccssbu-dar-equipment-ab}
\end{figure}
\FloatBarrier

\subsection{d-CCSSBU Charts for CAR of Equipment A and B Using Generated
jpeg
Files}\label{d-ccssbu-charts-for-car-of-equipment-a-and-b-using-generated-jpeg-files}

\begin{figure}[H]
\centering
\begin{tabular}{cc}
\includegraphics[width=0.48\textwidth]{d_CCSSBU_L_CAR_Equipment_A.jpeg} &
\includegraphics[width=0.48\textwidth]{d_CCSSBU_L_CAR_Equipment_B.jpeg} \\

\includegraphics[width=0.48\textwidth]{d_CCSSBU_S_CAR_Equipment_A.jpeg} &
\includegraphics[width=0.48\textwidth]{d_CCSSBU_S_CAR_Equipment_B.jpeg} \\
\end{tabular}
\caption{d-CCSSBU-L and d-CCSSBU-S charts for CAR of Equipment A and B.}
\label{fig:d-ccssbu-car-equipment-ab}
\end{figure}
\FloatBarrier

\subsection{d-CCSSBU Charts for Mic of Equipment A and B Using Generated
jpeg
Files}\label{d-ccssbu-charts-for-mic-of-equipment-a-and-b-using-generated-jpeg-files}

\begin{figure}[H]
\centering
\begin{tabular}{cc}
\includegraphics[width=0.48\textwidth]{d_CCSSBU_L_Mic_Equipment_A.jpeg} &
\includegraphics[width=0.48\textwidth]{d_CCSSBU_L_Mic_Equipment_B.jpeg} \\

\includegraphics[width=0.48\textwidth]{d_CCSSBU_S_Mic_Equipment_A.jpeg} &
\includegraphics[width=0.48\textwidth]{d_CCSSBU_S_Mic_Equipment_B.jpeg} \\
\end{tabular}
\caption{d-CCSSBU-L and d-CCSSBU-S charts for Mic of Equipment A and B.}
\label{fig:d-ccssbu-mic-equipment-ab}
\end{figure}
\FloatBarrier

\section{Poisson u Chart Illustration for Mic
Data}\label{poisson-u-chart-illustration-for-mic-data}

The following is three-sigma Poison u chart function
\texttt{poisson\_mic\_u\_chart()}.

The following is exact Poison u chart function
\texttt{poisson\_mic\_u\_chart\_exact()}.

\pandocbounded{\includegraphics[keepaspectratio]{multicluster_SPC_paper_code_files/figure-pdf/run-poisson-u-chart-exact-1.pdf}}

\subsection{Three Sigma Poisson u Chart for Mic Data of Equipment
A}\label{three-sigma-poisson-u-chart-for-mic-data-of-equipment-a}

\FloatBarrier

\begin{table}[!h]
\centering
\caption{\label{tab:display-poisson-u-decision-summary-equipment-a-v3}Mic data Poisson assessment and chart selection for Equipment A.\label{tab:poisson-u-assessment-equipment-a}}
\centering
\resizebox{\ifdim\width>\linewidth\linewidth\else\width\fi}{!}{
\fontsize{8}{10}\selectfont
\begin{tabular}[t]{llrrlrll}
\toprule
Equipment & Case & USL & GOF p & GOF pass & Dispersion & Overdispersion pass & Mic Poisson\\
\midrule
Equipment A & Case 3 & 25 & 0.627 & TRUE & 1 & TRUE & TRUE\\
\bottomrule
\end{tabular}}
\end{table}

\pandocbounded{\includegraphics[keepaspectratio]{multicluster_SPC_paper_code_files/figure-pdf/display-selected-poisson-u-chart-equipment-a-v31-1.pdf}}

\subsection{Three Sigma Poisson u Chart and Exact Poisson u Chart Using
Generated jpeg
Files}\label{three-sigma-poisson-u-chart-and-exact-poisson-u-chart-using-generated-jpeg-files}

\begin{figure}[H]
\centering
\begin{tabular}{cc}
\includegraphics[width=0.48\textwidth]{Poisson_U_Mic_Equipment_A.jpeg} &
\includegraphics[width=0.48\textwidth]{Poisson_U_Exact_Mic_Equipment_A.jpeg} 
\end{tabular}
\caption{Comparison of conventional and exact Poisson $u$ charts for Mic data from Equipment A.}
\label{fig:poisson-u-comparison-equipment-a}
\end{figure}
\FloatBarrier

\subsection{Non-Poisson Mic Data of Equipment
B}\label{non-poisson-mic-data-of-equipment-b}

\begin{table}[!h]
\centering
\caption{\label{tab:poisson-u-assessment-equipment-b}Mic data Poisson assessment and chart selection for Equipment B.}
\centering
\resizebox{\ifdim\width>\linewidth\linewidth\else\width\fi}{!}{
\fontsize{8}{10}\selectfont
\begin{tabular}[t]{llrrlrll}
\toprule
Equipment & Case & USL & GOF p & GOF pass & Dispersion & Overdispersion pass & Mic Poisson\\
\midrule
Equipment B & Case 3 & 25 & 0.0274 & FALSE & 1.6286 & FALSE & FALSE\\
\bottomrule
\end{tabular}}
\end{table}




\end{document}
